\chapter{不等式}


本章覆盖知识点如下：
\begin{introduction}
  \item 含参讨论
  \item 跨模块综合
  \item 新定义理解
\end{introduction}

\section{含参讨论}

参数与整数解、无解问题结合，需要精确分析端点。核心在于已知解集或整数解的情况，反求参数的取值范围，需要熟练掌握数轴分析和端点值的取舍。

\subsection{朱韬数学-不等式-1}

若数 \(a\) 使关于 \(x\) 的方程
\[
\frac{ax + 1}{2} = -\frac{7x}{3} - 1
\]
有非负数解，且关于 \(y\) 的不等式组
\[
\begin{cases} 
\frac{y - 1}{2} - 2 < \frac{7 - 2y}{2} \\ 
2y + 1 > a - 2y 
\end{cases}
\]

恰好有两个偶数解，则符合条件的所有整数 \(a\) 的和是\_\_\_\_\_\_。

（答案见第~\pageref{朱韬数学-不等式-1}~页）




\subsection{2024北京东城期末}

已知关于 \(x\) 的不等式组
\[
\begin{cases}
x \ge a, \\
2x - 3 < 1,
\end{cases}
\]
有三个整数解，则 \(a\) 的取值范围是 \_\_\_\_\_\_。

（答案见第~\pageref{2024北京东城期末}~页）

\subsection{2025北京二中期末}

对于实数 \(x_0, d\)（其中 \(d > 0\)），不等式 \(x_0 - d < x < x_0 + d\) 的解集构成“\(x_0\) 的 \(d\) 邻域”。

\begin{enumerate}
    \item 不等式 \(|x| + 3 < 2\) 的解集构成“\underline{\hspace{2cm}}邻域”；
    \item 不等式 \(A\) 的解集构成“\(a+1\) 的 \(2\) 邻域”，不等式 \(B\) 的解集构成“\(2\) 的 \(1\) 邻域”。若由 \(A, B\) 组成的不等式组的解集构成“\(\left(\dfrac{a}{2}+2\right)\) 的 \(\left(\dfrac{a}{2}+1\right)\) 邻域”，则 \(a\) 的取值范围是 \underline{\hspace{2cm}}。
\end{enumerate}

（答案见第~\pageref{2025北京二中期末}~页）

\subsection{2024北京房山期末}

甲、乙、丙三人做写数字的游戏，三个人写的数字要同时满足以下四个条件：

\begin{enumerate}
    \item 乙写的数字的一半大于甲写的数字；
    \item 丙写的数字不大于甲写的数字；
    \item 丙写的数字的 3 倍大于乙写的数字；
    \item 甲、乙、丙三人写的数字均为正整数。
\end{enumerate}

则三人所写数字之和的最小值为（\quad ）

A. 4 \quad B. 7 \quad C. 9 \quad D. 13

（答案见第~\pageref{2024北京房山期末}~页）



\subsection{不等式组含参问题-基础}

关于 \(x, y\) 的方程组 \(\begin{cases} x + 3y = 3 - 2k \\ 3x + y = 1 + k \end{cases}\) 的解满足 \(x + y > 0\)，且关于 \(x\) 的不等式组 \(\begin{cases} x - 2(x - 1) \leq 3 \\ \dfrac{2k + x}{3} \geq x \end{cases}\) 有解，则符合条件的整数 \(k\) 的值的和为（\quad ）

A. 2 \quad B. 3 \quad C. 4 \quad D. 5

\subsection{不等式组含参问题-基础2}

已知关于 \(x\) 的不等式 \((1+2a)x > 1\) 的解集为 \(x < \dfrac{1}{1+2a}\)，则 \(a\) 的取值范围是（\hspace{1cm}）

\noindent
\makebox[0.23\linewidth][l]{A. \(a < -\dfrac{1}{2}\)}%
\hfill
\makebox[0.23\linewidth][l]{B. \(a > -\dfrac{1}{2}\)}%
\hfill
\makebox[0.23\linewidth][l]{C. \(a < \dfrac{1}{2}\)}%
\hfill
\makebox[0.23\linewidth][l]{D. \(a > \dfrac{1}{2}\)}

（答案见第~\pageref{不等式组含参问题-基础2}~页）

\subsection{不等式含参问题}

已知关于 \(x, y\) 的方程组 \(\begin{cases} 2x - y = 1 - 3a \\ x + 3y = 2a + 4 \end{cases}\) 的解都为非负数，且满足 \(2a + b = 5\)，\(2 \leq b \leq 5\)，若 \(z = a - b\)，则 \(z\) 的取值范围是\_\_\_\_\_\_。



\section{不等式（组）与方程（组）的综合应用}

将方程、表格、不等式范围和新定义结合,压轴题的另一个方向是将二元一次方程组与不等式组结合，往往需要先用参数表示方程组的解，再根据解的“非负”“大小比较”或“整数”等条件列不等式。




\subsection{2025北京二中期末}

小明为了方便探究关于 \(x, y\) 的二元一次方程 \(x+y=3\) 解的规律，把 \(x\) 和 \(y\) 的部分值分别填入如表（\(x\) 的值从左到右依次增大）。

\[
\begin{array}{|c|c|c|c|c|c|}
\hline
x & -7 & -4 & 0 & 2 & 8 \\
\hline
y & 10 & 7 & p & 1 & -5 \\
\hline
\end{array}
\]

(1) \(p\) 的值为 \_\_\_\_\_\_。

(2) 下列方程中，与 \(x+y=3\) 组成方程组，在 \(-7<x<8\) 范围内有解的是 \_\_\_\_\_\_（填正确的序号）。
\begin{itemize}
\item ① \(3x-y=1\);
\item ② \(2x+y=-5\);
\item ③ \(x+2y=-4\)。
\end{itemize}

(3) 已知关于 \(x, y\) 的二元一次方程 \(cx+dy=1\)（\(c \neq 0, d \neq 0\)）的部分解如表所示：

\[
\begin{array}{|c|c|c|c|c|}
\hline
x & -7 & 2 & \cdots & 8 \\
\hline
y & -2 & n & \cdots & n+6 \\
\hline
\end{array}
\]

则方程组 \(\begin{cases} x+y=3 \\ cx+dy=1 \end{cases}\) 的解为 \_\_\_\_\_\_。



（答案见第~\pageref{2025北京二中期末}~页）


\subsection{2025北京门头沟期末}

暑期临近，某服装店老板计划购进甲、乙两种童装T恤，已知购进甲种T恤2件和乙种T恤3件共需310元；  
购进甲种T恤1件和乙种T恤2件共需190元。  

\begin{enumerate}
    \item 求甲、乙两种T恤每件的进价分别是多少元？  
    \item 为满足市场需求，服装店需购进甲、乙两种T恤共100件，要求购买两种T恤的总费用不超过6540元，并且购买的甲种T恤的数量的三倍不超过乙种T恤的数量，请你通过计算，确定服装店购买甲、乙两种T恤的购买方案。
\end{enumerate}

（答案见第~\pageref{2025北京门头沟期末}~页）


\section{新定义与阅读理解型不等式}

这类题目给出新的运算规则或概念，要求学生在理解的基础上应用不等式解决问题，非常考验综合能力。

\subsection{2025北京二中期末2}

对于实数 \(a, b\)，定义 \(\max\{a, b\}\)：当 \(a \geq b\) 时取 \(a\)，当 \(a < b\) 时取 \(b\)。

已知 \(\max\{2x + 1,\; 3x - 1\} = 3\)，求 \(x\) 的值。

（答案见第~\pageref{2025北京二中期末2}~页）


\subsection{新定义-集合思想}

将二元一次方程组的解中的所有数的全体记为 \(M\)，将不等式（组）的解集记为 \(N\)，给出定义：若 \(M\) 中的数都在 \(N\) 内，则称 \(M\) 被 \(N\) 包含；若 \(M\) 中至少有一个数不在 \(N\) 内，则称 \(M\) 不能被 \(N\) 包含。

如：方程组 \(\begin{cases} x = 0 \\ x + y = 2 \end{cases}\) 的解为 \(\begin{cases} x = 0 \\ y = 2 \end{cases}\)，记 \(A: (0, 2)\)；方程组 \(\begin{cases} x = 0 \\ x + y = 4 \end{cases}\) 的解为 \(\begin{cases} x = 0 \\ y = 4 \end{cases}\)，记 \(B: (0, 4)\)；不等式 \(x - 3 < 0\) 的解集为 \(x < 3\)，记 \(H: x < 3\)。因为 \(0, 2\) 都在 \(H\) 内，所以 \(A\) 被 \(H\) 包含；因为 \(4\) 不在 \(H\) 内，所以 \(B\) 不能被 \(H\) 包含。

(1) 将方程组 \(\begin{cases} 2x - y = 5 \\ 3x + 4y = 2 \end{cases}\) 的解中的所有数的全体记为 \(C\)，将不等式 \(x + 1 \geq 0\) 的解集记为 \(D\)，请问 \(C\) 能否被 \(D\) 包含？说明理由。

(2) 将关于 \(x, y\) 的方程组 \(\begin{cases} 2x + 3y - 5a = -1 \\ x - 2y + a = 3 \end{cases}\) 的解中的所有数的全体记为 \(E\)，将不等式组 \(\begin{cases} 3(x - 2) \geq x - 4 \\ 2x + 1 \geq x - 1 \end{cases}\) 的解集记为 \(F\)，若 \(E\) 不能被 \(F\) 包含，求实数 \(a\) 的取值范围。


（答案见第~\pageref{新定义-集合思想}~页）


\subsection{二中教育集团期末第26题}

已知二元一次方程组 \(M\) 的解为 $\begin{cases} x = a \\ y = b \end{cases}$，在数轴上实数 \(a\) 所对的点为 \(A\)，实数 \(b\) 所对的点为 \(B\)。若在线段 \(AB\) 上存在 \(K\) 个整数，则称二元一次方程组 \(M\) 为 \(K$ 系方程组。

1. 二元一次方程组 $\begin{cases} x + y = 4 \\ x = 1.5 \end{cases}$ 是 \_\_\_\_\_\_ 系方程组。

2. 关于 \(x, y\) 的二元一次方程组 $\begin{cases} x + y = 2.5 + 2m \\ x - 3y = 2.5 - 6m \end{cases}$ 是 3 系方程组，直接写出 \(m\) 的取值范围。

3. 关于 \(x, y\) 的二元一次方程组 $\begin{cases} x - y = 2n - 1.6 \\ x = n + 0.4 \end{cases}$ 是 2 系方程组，直接写出 \(n\) 的取值范围。

（答案见第~\pageref{二中教育集团期末第26题}~页）



\section{不等式与最值综合}

\subsection{不等式与最值综合-题目1}

已知实数 \(x\) 满足 \(\begin{cases} 5(x+1) \geq 3x-1 \\ \frac{1}{2}x-1 \leq 7-\frac{3}{2}x \end{cases}\)，若 \(S = |x-1|+|x+1|\) 的最大值为 \(m\)，最小值为 \(n\)，则 \(mn = \underline{\quad }\)。


(答案见第~\pageref{不等式与最值综合-题目1}~页)


\subsection{不等式与坐标系-题目1}

如果关于 \(x\) 的不等式组 
\(
\begin{cases} 
2x - a \geq 0 \\ 
b - 3x \geq 0 
\end{cases}
\)
仅有四个整数解为 \(-2, -1, 0, 1\)，若 \(P(a, b)\) 在第二象限，那么满足上述条件的整数 \(a, b\) 组成的点 \(P\) 的坐标有（ \quad）

A. 2 个 \quad B. 4 个 \quad C. 6 个 \quad D. 9 个

(答案见第~\pageref{不等式与坐标系-题目1}~页)

